\chapter[Second-order modulation]{A simulation of second-order amplitude and phase modulation}
\label{ch:second-order-modulation}
Phase--amplitude coupling (PAC) is usually stated as a first-order
relation: the phase of a slower rhythm modulates the envelope of a faster
one \cite{canolty2006,tort2010}. Hierarchical, nested accounts of LFP
rhythms \cite{hyafil2015} invite a further question: is the
\emph{coupling coefficient itself} modulated by a third, still slower
phase? This chapter poses a generative model in which the answer is known,
derives what an estimator should return, and records what the estimator
actually returns once filtering, noise and finite data intervene. It uses
Bedrosian's theorem (Chapter~\ref{ch:classical-sampling}) and the Hilbert
multiplier identity of Chapter~5 at every level of the hierarchy.

The code is \path{simulate_lfp_modulation.m} and
\path{second_order_modulation.m}, with the regression test
\path{matlab/tests/test_second_order_modulation.m}.

\section{Orders of modulation}
\begin{definition}[Order of a modulation]
Let $x(t)=\sum_k A_k(t)\cos\phi_k(t)$ be a sum of narrowband components.
A parameter of component $k$ (its envelope $A_k$ or the deviation of
$\phi_k$ from a linear ramp) is \emph{modulated to first order} by
component $j$ if it is a nonconstant function of $\phi_j$ with fixed
coefficients. It is modulated \emph{to second order} by a pair $(j,l)$
if it depends on $\phi_j$ through a coefficient that is itself a
nonconstant function of $\phi_l$.
\end{definition}
The word ``order'' here counts levels of a modulation hierarchy, not the
order of a moment or of a filter.

\section{The generative model}
Write $s,m,f$ for slow, middle and fast components (defaults
$f_s=1.5$, $f_m=8$, $f_f=60$\,Hz: delta, theta, gamma). With independent
standard Wiener processes $W_\cdot$, phase-diffusion coefficients
$\sigma_\cdot$ (rad$/\sqrt{\mathrm s}$) and uniform initial phases
$u_\cdot$,
\begin{align*}
\phi_s&=2\pi f_st+u_s+\sigma_sW_s,\\
A_m&=a_m\bigl(1+m_m\cos(\phi_s-\psi_m)\bigr), &
\phi_m&=2\pi f_mt+u_m+\sigma_mW_m+\beta_m\sin(\phi_s-\chi_m),\\
m_f&=m_0\bigl(1+\kappa\cos(\phi_s-\psi_2)\bigr), &
b_f&=b_0\bigl(1+\lambda\cos(\phi_s-\chi_2)\bigr),\\
A_f&=a_f\bigl(1+m_f\cos(\phi_m-\psi_f)\bigr), &
\phi_f&=2\pi f_ft+u_f+\sigma_fW_f+b_f\sin(\phi_m-\chi_f),
\end{align*}
and $x=a_s\cos\phi_s+A_m\cos\phi_m+A_f\cos\phi_f+\xi$, where $\xi$ is a
$1/f^{\chi}$ Gaussian background plus white measurement noise. The middle
component is amplitude- and phase-modulated by the slow phase (first
order); the fast component is amplitude- and phase-modulated by the middle
phase with depth $m_f$ and index $b_f$, and these are modulated by the slow
phase (second order). Setting $\kappa=\lambda=0$ recovers ordinary
first-order PAC and phase--phase modulation. Envelopes are nonnegative
when $m_m\le1$ and $m_0(1+\kappa)\le1$, which the generator enforces.

Phase diffusion is not decoration. With $f_m/f_s=16/3$ the deterministic
phases would be commensurate, $\phi_m$ would take only a few values for
each $\phi_s$, and the conditional expectations below would not exist in
the form stated. Diffusion of the relative phase
$\phi_m-\tfrac{16}{3}\phi_s$, with variance rate
$\sigma_m^2+(16/3)^2\sigma_s^2\approx5.2$\,rad$^2$/s at the defaults,
makes $\phi_m$ asymptotically uniform given $\phi_s$.

\section{Spectral signature}
\begin{proposition}[Lines produced by second-order AM]
\label{prop:som-am-lines}
With $b_f=0$, the fast term expands as
\begin{align*}
A_f\cos\phi_f={}&a_f\cos\phi_f
+\tfrac{a_fm_0}{2}\sum_{\pm}\cos\bigl(\phi_f\pm(\phi_m-\psi_f)\bigr)\\
&+\tfrac{a_fm_0\kappa}{4}\sum_{\pm,\pm}
\cos\bigl(\phi_f\pm(\phi_m-\psi_f)\pm(\phi_s-\psi_2)\bigr).
\end{align*}
Hence second-order AM places four lines at $f_f\pm f_m\pm f_s$, each of
amplitude $a_fm_0\kappa/4$, whose phases are locked to
$\phi_f\pm\phi_m\pm\phi_s$.
\end{proposition}
\begin{proof}
Expand $A_f=a_f+a_fm_0\cos(\phi_m-\psi_f)
+a_fm_0\kappa\cos(\phi_m-\psi_f)\cos(\phi_s-\psi_2)$ and apply
$\cos\alpha\cos\beta=\tfrac12[\cos(\alpha+\beta)+\cos(\alpha-\beta)]$
once to the second term and twice to the third.
\end{proof}

\begin{proposition}[Small-index PM lines]
With $m_f=0$, $u=\phi_m-\chi_f$ and $v=\phi_s-\chi_2$, the Jacobi--Anger
expansion $\cos(\theta+\beta\sin u)=\sum_{n\in\Z}J_n(\beta)\cos(\theta+nu)$
and $J_{-1}=-J_1$, $J_1(\beta)=\beta/2+O(\beta^3)$ give, to first order in
$b_0$,
\[
\cos\phi_f\approx\cos\theta
+\tfrac{b_0}{2}\bigl[\cos(\theta+u)-\cos(\theta-u)\bigr]
+\tfrac{b_0\lambda}{4}\sum_{\pm}\bigl[\cos(\theta+u\pm v)-\cos(\theta-u\pm v)\bigr],
\]
with $\theta$ the unmodulated fast phase.
\end{proposition}
Second-order AM and second-order PM therefore occupy the \emph{same}
frequencies $f_f\pm f_m\pm f_s$. They differ in the sign relation between
upper and lower sidebands (equal for AM, opposite for PM), which is the
classical AM/PM quadrature distinction. A power spectrum, which discards
these phases, cannot tell them apart; nor can it tell either from
independent oscillations that happen to sit at those frequencies.

The phase relation $\phi_{f_f+f_m+f_s}-\phi_f-\phi_m-\phi_s=\text{const}$
involves four components, so it is a trispectral (fourth-order)
statistic. A bispectrum can respond only through a triad. When the middle
component is itself AM- or PM-modulated by $\phi_s$, it carries sidebands
at $f_m\pm f_s$, the triad $(f_m+f_s,f_f,f_f+f_m+f_s)$ exists, and the
bispectrum can respond; it then cannot separate first-order modulation of
the middle rhythm from second-order modulation of the fast one.

\section{Spectral separation at every level}
Bedrosian's theorem (Chapter~\ref{ch:classical-sampling}) gives a
sufficient condition for the analytic signal of $A\cos\phi$ to return
$A$ and $\phi$: the envelope spectrum lies inside $(-B,B)$ and the
carrier spectrum outside it. With phase modulation the carrier is not a
line. Carson's rule \cite{carson1922} approximates the band occupied by a
sinusoidally phase-modulated carrier of index $\beta$ and modulating
frequency $f_{\mathrm{mod}}$ as $2(\beta+1)f_{\mathrm{mod}}$; it is an
engineering approximation (it retains roughly $98\%$ of the power), not a
theorem. The resulting conditions, used by the generator, are
\begin{align*}
\text{middle:}&\quad f_s<f_m-(\beta_m+1)f_s,\\
\text{fast:}&\quad f_m+f_s<f_f-(b_0(1+\lambda)+1)f_m,
\end{align*}
and an analysis passband of half-width
\[
h_m=c\bigl[(\beta_m+1)f_s+f_s\bigr],\qquad
h_f=c\bigl[(b_0(1+\lambda)+1)f_m+f_m+f_s\bigr],
\]
with a margin $c\ge1$ (default $1.2$) against Butterworth roll-off. At the
defaults this gives $[2.96,13.04]$\,Hz and $[30.36,89.64]$\,Hz, and a slow
band $f_s\pm cf_s/2=[0.6,2.4]$\,Hz. The generator returns these bands in
\texttt{recommended\_band\_hz} and flags every violated condition in
\texttt{constraints}; a band chosen by convention (``theta $5$--$11$\,Hz'')
rather than by the model is the source of two of the failures in
Section~\ref{sec:som-failures}.

\section{Estimator targets}
\begin{proposition}[Conditional mean-vector length]
\label{prop:som-mvl}
Assume $\phi_m$ is uniform on the circle conditionally on $\phi_s$ and
independent of the diffusion of $\phi_f$. Then
\[
\frac{\E[A_fe^{i\phi_m}\mid\phi_s]}{\E[A_f\mid\phi_s]}
=\frac{m_f(\phi_s)}{2}e^{i\psi_f}.
\]
In particular, the normalised mean-vector length equals $m_0/2$
marginally and $m_f(\phi_s)/2$ conditionally.
\end{proposition}
\begin{proof}
Conditionally on $\phi_s$, $m_f$ is fixed. Uniformity gives
$\E[e^{i\phi_m}]=0$ and $\E[\cos(\phi_m-\psi_f)e^{i\phi_m}]
=\tfrac12\E[e^{i\psi_f}+e^{i(2\phi_m-\psi_f)}]=\tfrac12e^{i\psi_f}$.
Hence the numerator is $a_fm_fe^{i\psi_f}/2$ and the denominator is
$a_f(1+m_f\E\cos(\phi_m-\psi_f))=a_f$.
\end{proof}

\begin{proposition}[Conditional PM index]
Under the same assumption, with
$d_f=b_f\sin(\phi_m-\chi_f)$ the phase deviation,
$2\,\lvert\E[d_fe^{i\phi_m}\mid\phi_s]\rvert=b_f(\phi_s)$.
\end{proposition}
\begin{proof}
$\sin(\phi_m-\chi_f)e^{i\phi_m}
=\tfrac1{2i}\bigl(e^{i(2\phi_m-\chi_f)}-e^{i\chi_f}\bigr)$; the first
term has zero conditional mean, so the expectation is
$-b_fe^{i\chi_f}/(2i)$, of modulus $b_f/2$.
\end{proof}

\paragraph{Centring.} In a finite bin the empirical mean of $e^{i\phi_m}$
is not zero. Writing $\bar{\,\cdot\,}$ for a bin average,
$\overline{Ae^{i\phi}}=\overline{(A-\bar A)e^{i\phi}}
+\bar A\,\overline{e^{i\phi}}$: the second term leaks the mean envelope
into the coupling estimate. In the default simulation, with $200$\,s and
$12$ bins, $\lvert\overline{e^{i\phi_m}}\rvert\approx0.01$--$0.04$ per bin,
which on the \emph{exact} phases produced per-bin errors of up to $0.09$
in $\hat m_f$. The estimator therefore uses the centred projection
$\overline{(A-\bar A)e^{i\phi}}$, which has the same target under
uniformity and is the phase-clustering correction of
\cite{vandriel2015}; the worst-bin error on exact phases fell to
$0.024$, attributable to sampling.

\begin{lemma}[Binning attenuates the first harmonic]
If $y(\theta)=c_0+\delta\cos(\theta-\psi)$ is averaged over bins of width
$w=2\pi/M$, the bin means are $c_0+\delta\,\mathrm{sinc}_M\cos(c_b-\psi)$
at bin centres $c_b$, with $\mathrm{sinc}_M=\sin(\pi/M)/(\pi/M)$.
\end{lemma}
\begin{proof}
$\frac1w\int_{c-w/2}^{c+w/2}\cos(\theta-\psi)\,d\theta
=\frac{2\sin(w/2)}{w}\cos(c-\psi)$.
\end{proof}
For $M=12$ the factor is $0.9886$; the estimator divides the fitted
first-harmonic amplitude by it. This assumes $\phi_s$ is roughly uniform
within each bin.

\paragraph{The estimator.} After zero-phase band-pass filtering and Hilbert
transformation of each band, and removal of filter edges, the slow phase
is divided into $M$ bins. In each bin
\[
\hat m(b)=\frac{2\lvert\overline{(\hat A_f-\bar{\hat A}_f)e^{i\hat\phi_m}}\rvert}
{\bar{\hat A}_f},\qquad
\hat\beta(b)=2\lvert\overline{(\hat d_f-\bar{\hat d}_f)e^{i\hat\phi_m}}\rvert,
\]
where $\hat d_f$ is the unwrapped fast phase after a zero-phase high-pass
below the middle band, which removes the carrier ramp and the diffusion
drift. Each profile is fitted by least squares to
$c_0+c\cos\theta+s\sin\theta$; the code reports $c_0$, the corrected
absolute second-order term $\Delta=\lvert c+is\rvert/\mathrm{sinc}_M$, the
ratio $\Delta/c_0$ and the phase $\arg(c+is)$. For the model these
estimate $(m_0,m_0\kappa,\kappa,\psi_2)$ and $(b_0,b_0\lambda,\lambda,\chi_2)$.

\begin{caution}
Report $\Delta$, not only the ratio. When first-order coupling is near
zero, $\Delta/c_0$ is a ratio of two noise-dominated quantities; in a
simulation with no coupling at all it averaged $0.21$ with standard
deviation $0.18$ across six seeds, while $c_0=0.008$: a meaningless number.
\end{caution}

\paragraph{A null for the second order.} Circularly shifting $\hat\phi_s$
by at least $10$\,s, while leaving $(\hat\phi_m,\hat A_f,\hat d_f)$ in
place, preserves first-order coupling and the circular within-series
structure. The reference distribution is appropriate only when admissible
relative shifts are exchangeable under the chosen null and disrupt the
dependence measured by $\Delta$. It does not universally encode ``no
dependence of the coupling on $\phi_s$''. For a periodic slow phase,
$\phi_s(t-\tau)=\phi_s(t)-\omega_s\tau$ modulo $2\pi$, a shift can simply
rotate the coupling profile: its preferred phase changes while its
first-harmonic amplitude stays unchanged. The test can then lack power
even when second-order modulation is present. A nonsignificant result
does not establish absence of dependence; phase diffusion and the shift
range must be assessed for each recording. The zero-phase-diffusion
regression in \path{matlab/tests/test_second_order_modulation.m} retains
this failure case explicitly. In $40$
null realisations ($\kappa=\lambda=0$, $120$\,s, $99$ shifts) the
nominal-$5\%$ test rejected $0/40$ times for AM and $3/40$ for PM
($P(\ge3)\approx0.32$ under a $5\%$ rate); the NumPy mirror gave $0/40$
and $4/40$. Both are compatible with the nominal level, which more runs
should confirm.

\section{What the simulation returned}
Every number in this section and the next is regenerated by
\path{examples/matlab/second_order_modulation_study.m}, which runs the
MATLAB functions over fixed seeds and writes a results file that replaces
the tables below when present in \path{figures/generated/}. Seeds and
random generators differ between runs on different platforms, so values
agree in distribution rather than digit by digit.

% Generated tables (study output) take precedence over the fallback.
\InputIfFileExists{figures/generated/second_order_results.tex}{}{%
Fallback (no generated file): the study run on 20 October 2026 with the
MATLAB sources executed under GNU Octave~8.4 and the Signal package, with
test-only compatibility shims for \texttt{RandStream}, \texttt{discretize}
and second-order-section \texttt{filtfilt}. On identical input the MATLAB
estimator and an independent NumPy implementation agreed to four
decimals; a MATLAB run is still pending.
\begin{center}
\begin{tabular}{lrrr}
\toprule
Quantity & True & Mean & SD\\
\midrule
$m_0$ & $0.500$ & $0.481$ & $0.005$\\
$m_0\kappa$ & $0.300$ & $0.284$ & $0.004$\\
$\kappa$ & $0.600$ & $0.589$ & $0.005$\\
$\psi_2$ (rad) & $1.047$ & $1.056$ & $0.011$\\
$b_0$ & $0.600$ & $0.595$ & $0.003$\\
$b_0\lambda$ & $0.300$ & $0.291$ & $0.009$\\
$\lambda$ & $0.500$ & $0.490$ & $0.016$\\
$\chi_2$ (rad) & $-0.785$ & $-0.739$ & $0.037$\\
normalised MVL & $0.250$ & $0.241$ & $0.002$\\
\bottomrule
\end{tabular}
\end{center}
\noindent GNU Octave 8.4.0, 8 seeds of $200$\,s. Null calibration: 40 realisations of $120$\,s, 99 shifts; nominal-$5\%$ rejections 0 (AM) and 3 (PM); $P(X\ge3)=0.32$ under a $5\%$ rate.

\begin{center}
\small
\begin{tabular}{lrrrrc}
\toprule
Scenario & $\hat c_0^{\mathrm{AM}}$ & $\hat\Delta_{\mathrm{AM}}$ & $\hat\Delta_{\mathrm{AM}}/\hat c_0$ & $\hat\Delta_{\mathrm{PM}}$ & sig.\ AM/PM\\
\midrule
No coupling & $0.008$ & $0.0016$ & $0.21\pm0.18$ & $0.0057$ & --\\
Harmonics, slow AM of middle & $0.135$ & $0.1141$ & $0.84\pm0.05$ & $0.0759$ & 6/6, 6/6\\
Harmonics, no slow AM & $0.120$ & $0.0248$ & $0.21\pm0.05$ & $0.0178$ & --\\
$\kappa=0$ ($\lambda=0.5$), $35$--$85$\,Hz, $\beta_m=0.8$ & $0.481$ & $0.0069$ & $0.01\pm0.00$ & $0.2903$ & 1/6, 6/6\\
$\kappa=0$, $35$--$85$\,Hz, $\beta_m=0$ & $0.483$ & $0.0037$ & $0.01\pm0.00$ & $0.2871$ & 0/6, 6/6\\
$\kappa=0$, $25$--$95$\,Hz, $\beta_m=0.8$ & $0.479$ & $0.0036$ & $0.01\pm0.00$ & $0.2924$ & 0/6, 6/6\\
\bottomrule
\end{tabular}
\end{center}
\noindent Middle AM depth (truth $0.4$): $0.390$ ($\beta_m=0$, $5$--$11$\,Hz), $0.356$ ($\beta_m=0.8$, $5$--$11$\,Hz), $0.385$ (Carson band). Instantaneous-frequency PM estimator ($35$--$85$\,Hz): ratio $0.583$, phase $-0.624$ with $\beta_m=0.8$ (prediction $0.615$, $-0.612$); $0.480$, $-0.784$ with $\beta_m=0$ (truth $0.5$, $-0.785$). $\hat\chi_2$: $-0.758$ ($\psi_2=\pi/3$), $-0.807$ ($\kappa=0$), $-0.844$ ($\psi_2=-3\pi/4$), $-0.793$ (no background noise).
}

% Optional generated figure from the same study.
\IfFileExists{figures/generated/fig_second_order_recovery.pdf}{%
\begin{figure}[t]
  \centering
  \includegraphics[width=\textwidth]{figures/generated/fig_second_order_recovery.pdf}
  \caption{One simulated realisation. Left: one second of the LFP.
  Centre and right: estimated PAC depth and PM index in twelve slow-phase
  bins against the true $m_f(\phi_s)$ and $b_f(\phi_s)$.}
  \label{fig:second-order-recovery}
\end{figure}
}{}

The depths are biased downward by a few per cent: the outer sidebands
$f_f\pm f_m\pm f_s$ and the middle band's PM spread fall where a
fourth-order Butterworth filter applied forward and backward no longer has
unit gain. The bias in $\chi_2$ is a separate mechanism, discussed next.
With $\kappa=\lambda=0$ the absolute terms averaged $0.005$ (AM) and
$0.010$ (PM) over the $40$ null realisations, the size of the upward bias
of a modulus of noisy means. In the $\kappa=0$ rows of the confound table
$\lambda$ stays at $0.5$, so the PM detections there are correct.

\section{Four ways to obtain second-order modulation without it}
\label{sec:som-failures}
Each of the following produced a slow-phase dependence of measured
coupling that the circular-shift null declares significant, in a model
where the second-order mechanism is absent or different. All were
isolated by switching one parameter at a time.

\begin{caution}[Non-sinusoidal middle rhythm]
Add harmonics $h_k\cos(k\phi_m)$, $h_k=0.6/k$, $k=2,\dots,10$, to the
middle rhythm and remove every true coupling. The harmonics of an $8$\,Hz
rhythm fall inside the fast band, so the fast envelope and phase become
periodic in $\phi_m$: spurious first-order PAC (normalised MVL $0.07$,
$\hat m_0\approx0.13$) and spurious first-order PM ($\hat b_0\approx0.10$). Because
the harmonics scale with $A_m$, which the slow phase modulates, the
spurious coupling also depends on $\phi_s$: $\hat\Delta_{\mathrm{AM}}\approx0.11$,
$\hat\Delta_{\mathrm{PM}}\approx0.08$, significant in $6/6$ seeds for both. Without
slow AM of the middle rhythm, $\hat\Delta_{\mathrm{AM}}$ fell from $0.114$ to
$0.025$ and the ratio from $0.84$ to $0.21$. Waveform shape
is the standard PAC confound \cite{aru2015,cole2019}; here it propagates
one level up. The generator flags it
(\texttt{harmonics\_outside\_fast\_band} false); the surrogate does not.
\end{caution}

\begin{caution}[Modulator frequency modulation and filter roll-off]
When the middle rhythm is phase-modulated by $\phi_s$, its instantaneous
frequency $f_m+\beta_mf_s\cos(\phi_s-\chi_m)$ moves with $\phi_s$, and so
do the fast sidebands $f_f\pm f_m$. If the fast band is narrower than the
Carson rule requires, the sidebands slide across the filter skirt and the
measured depth depends on $\phi_s$. With a conventional $35$--$85$\,Hz band
and $\kappa=0$: $\hat\Delta_{\mathrm{AM}}=0.0069\pm0.0019$, significant in
$1/6$ seeds ($2/6$ in the NumPy mirror); with $\beta_m=0$, $0.0037$ and
$0/6$; with a $25$--$95$\,Hz band, $0.0036$ and $0/6$. The bias is
consistent (about twice the baseline) even where single runs do not reach
significance. The effect is small next to a true $m_0\kappa=0.3$ and is
expected to reach significance more often in longer recordings, because
the surrogate spread shrinks with length while the bias does not. (An earlier, uncentred version of the estimator gave
$0.008$ and $4/6$.) The same
mechanism (PM-to-AM conversion by a non-flat passband) lowered the
first-order estimate of $m_m$, from a harmonic fit of the middle envelope
on $\hat\phi_s$, from $0.390$ ($\beta_m=0$) to $0.356$ ($\beta_m=0.8$) with a
$5$--$11$\,Hz middle band; the Carson band restored $0.385$.
\end{caution}

\begin{caution}[Instantaneous-frequency PM estimators]
An apparently natural estimator regresses the fast instantaneous frequency
on $\cos\hat\phi_m$. Differentiating $b_f\sin(\phi_m-\chi_f)$ gives
$b_f\dot\phi_m\cos(\phi_m-\chi_f)+\dot b_f\sin(\phi_m-\chi_f)$; the
estimated index is $\propto b_f\dot\phi_m$, and $\dot\phi_m$ is itself
modulated by $\phi_s$ when $\beta_m\ne0$. The measured profile is then
\[
(1+\lambda\cos(\phi_s-\chi_2))
\bigl(1+\beta_m\tfrac{f_s}{f_m}\cos(\phi_s-\chi_m)\bigr)
\approx1+\Re\bigl[(\lambda e^{-i\chi_2}+\beta_m\tfrac{f_s}{f_m}e^{-i\chi_m})
e^{i\phi_s}\bigr],
\]
which at the defaults predicts a ratio $0.615$ and phase $-0.61$\,rad
instead of $0.5$ and $-0.785$; the simulation (with a $35$--$85$\,Hz band)
returned $0.58$ and $-0.62$, and $0.48$ and $-0.78$ with $\beta_m=0$.
The phase-deviation estimator
used in the code has no $\dot\phi_m$ factor and is immune.
\end{caution}

\begin{caution}[AM leaking into PM through phase noise]
The phase of a noisy analytic signal is least reliable where its envelope
is small, and second-order AM makes the fast envelope's troughs deepest at
$\phi_s=\psi_2$. The PM estimate inherits a component at that slow phase:
$\hat\chi_2$ was $-0.758$ with $\psi_2=\pi/3$, $-0.807$ with $\kappa=0$,
and $-0.844$ with $\psi_2=-3\pi/4$ (truth $-0.785$, between-seed SD about
$0.04$); without background noise it was $-0.793$. The NumPy mirror gave
$-0.705$, $-0.774$, $-0.835$ and $-0.776$: the same ordering, and a shift
toward $\psi_2$ that the noiseless case removes. Second-order AM and PM estimates are therefore not
independent, and a significant PM result in the presence of strong AM
needs a simulation of this kind at matched signal-to-noise ratio.
\end{caution}

\section{What the model does and does not establish}
Recovery within a few per cent shows that the estimator approximates its
target under the model's assumptions; it is a statement about numerical
behaviour, not about the brain. The four failures above are each
sufficient to produce surrogate-significant ``second-order coupling''
without a second-order mechanism, and real LFPs contain all of their
ingredients: non-sinusoidal theta, theta frequency that varies with slower
state, conventional band choices and envelope-dependent noise. A claim of
second-order modulation in data therefore requires at least: bands
justified by the separation conditions above, waveform-shape controls,
the phase-deviation (not instantaneous-frequency) PM estimator, the
absolute term $\Delta$ with a surrogate reference justified by the observed
phase dynamics (including an assessment of circular-shift power), and a matched
simulation showing that the observed $\Delta$ exceeds what the confounds
produce at the data's signal-to-noise ratio.

\section*{Exercises for this chapter}
\begin{exercise}
Extend Proposition~\ref{prop:som-am-lines} to include $b_f\ne0$ and find
the lowest order in $(m_0,b_0)$ at which AM and PM sidebands interact.
\end{exercise}
\begin{exercise}
Show that the trispectrum at $(f_s,f_m,f_f)$ of the model with
$m_m=\beta_m=0$ is nonzero if and only if $\kappa\ne0$ or $\lambda\ne0$.
Which product of moments must be removed to make this a cumulant
statement?
\end{exercise}
\begin{exercise}
Compute the variance of $\lvert\overline{e^{i\phi}}\rvert^2$ for $N$
independent uniform phases and use it to explain the per-bin errors of the
uncentred estimator.
\end{exercise}
\begin{exercise}
Using \path{simulate_lfp_modulation.m}, find the smallest $f_f$ for which
the default fast band no longer satisfies the Bedrosian condition, and
measure how $\hat\Delta_{\mathrm{AM}}$ degrades below it.
\end{exercise}
\begin{exercise}
Design a waveform-shape control for the first caution: for example,
estimate coupling after removing the part of the fast band predicted by a
harmonic regression on $\hat\phi_m$, and test it on the harmonic
simulation.
\end{exercise}
